\section{广义黏滞方程的数学架构}

正文第二章给出了广义黏滞系数 $\eta = f(R,T,G)$，并在第 2.5 节末写出了指数形式 $\eta=\eta_0\exp(\alpha R+\beta T+\gamma G)$ 与三个系数各自的含义；那里没有展开的，是这个形式为什么取指数、三项为什么可乘、以及它落到可检验命题上还缺什么。这一节补的正是这三件事，同时说明这个形式的地位：它是一组待检验的假设形式，不是已经被确证的关系。

\subsection{内禀黏滞指数法则}

假设在固定任务与测量口径下，三类约束对对数转移代价的边际影响近似恒定，且暂不计入交互项，那么总黏滞系数可以写成

\begin{equation}
\eta = \eta_0 \exp\left(\alpha R + \beta T + \gamma G\right)
\end{equation}

式中 $\eta_0$ 是基准黏滞系数，取值取决于所讨论的环节本身，与三类约束无关；$R$、$T$、$G$ 分别是资源势垒、时间刚性和拓扑耦合的归一化度量，量纲为一；$\alpha$、$\beta$、$\gamma$ 是三个常数，表示各类约束的相对权重，其数值需由实证确定，预期均为正。指数形式的来源是：如果每一类约束按比例降低通过速率，而多类约束同时作用时降低倍率相乘，则总倍率的对数等于各类约束的线性叠加。这个形式与阿伦尼乌斯方程中速率常数对活化能的指数依赖属于同一族，其共同点是处理“需要通过一个障碍”的过程，而不是处理“在平坦空间中移动”的过程。

指数中的可加性，等价于相对代价倍率可乘。它要求$\partial\ln\eta/\partial R=\alpha$等边际关系在考察区间内近似不变。物理串联过程的历时可以相加，稳态管线的吞吐也可能由最慢环节限制，不能仅凭“串联”推出倍率相乘。出现交互作用时，可检验$\delta RT$等扩展项；采用哪种形式，要比较解释与样本外表现。

测量先固定一个结果口径。$\eta=C/C_{\mathrm{ref}}$中的$C$可以取费用、总工时或单位合格结果的平均通过时间，$C_{\mathrm{ref}}>0$取同任务的公开参照。成功率与代价的方向相反，若采用成功率作代理，需要先规定相应变换。测量基准确定以后，$\alpha$才可读为资源约束增加一单位时对数代价的边际变化。不同研究的变量、边界和参照一致，系数才有可比性。

三个自变量需要各自公开定义、基准与取值方式。归一化可以采用同类需求与可用能力的比值，或采用可重复的连续指标；0到5级的等级表则先用于序数判断。自变量应尽量由转移前的条件或独立资料构造，避免把待解释的费用与历时换名后再次作为解释变量。跨研究合并前先核对任务、刻度和边界，再比较系数与残差。

这个形式可以提出三类检查：残差是否仍按任务或介质呈现结构，参数能否在可比样本中复现，交互与阈值形式能否更好地解释观察。指数函数的边际支出平滑增长，并不自动产生相变或突变阈值。模型形状与具体临界机制需要分别建立。

等黏滞面上的置换比可以写成，

\begin{equation}
\left(\frac{\partial R}{\partial T}\right)_{\eta}=-\frac{\beta}{\alpha}
\end{equation}

式中，$R$ 为资源势垒的归一化度量，$T$ 为时间刚性的归一化度量，两者量纲为一；$\alpha$ 与 $\beta$ 分别为 $\ln\eta$ 对 $R$ 与对 $T$ 的偏导数，量纲为一；下标 $\eta$ 表示在黏滞系数保持不变的条件下取值。这个式子的含义是：要让黏滞保持不变，资源约束每放松一个单位，时间约束沿等值面收紧$\alpha/\beta$个单位，反之亦然，这正是正文所说的三组约束可以相互置换的定量形式，即压缩时间要投入更多资源、节省资源要拉长时间。适用条件必须写明：$\alpha$、$\beta$ 的数值须由实证确定，本式只给出置换比的结构；它成立的前提是可乘假设；一旦这个前提失效，就需要在指数上补一个交叉项，置换比随之不再是一个常数。

三条含义中，第一条最容易被误读。残差没有结构是模型成立的必要条件，不是充分条件：一个错误的模型也可能在小样本上给出无结构的残差。因此这条检验必须以介质为分组单元来做，而不是把不同介质的数据混在一起拟合；混在一起时，介质的差异会被吸收进残差，看起来像噪声，实际上是被忽略的变量。这一点决定了附录第三节的检验路线为什么必须按介质分别设计。

参数的跨期稳定性需要样本外检查。漂移可能来自技术变化、边界改变、变量定义或遗漏条件，不能据此直接判定因果方向反转。应保留定义与版本，比较替代解释，并说明模型在什么区间仍有预测能力。

还有一点关于自变量的写法需要交代。三类约束各自内部可能包含若干不同的成分，例如资源约束里既有材料的可得性也有资金的可得性，而公式里只允许出现一个数。把多维的约束压成一个标量，意味着研究者要先确定用什么方式汇总，例如取各项的最大值，或者按某个权重求和，而这个汇总方式与三个系数一样需要写明。汇总方式选得不好，误差会进入残差并被当成噪声；如果选得好，残差应当回到随机水平。这条要求不增加新的假设，它只是承认一件本来就必须回答的事：任何一个只允许输入一个数的形式，都要先交代这个数是怎么来的。

\subsection{社会技术德博拉数}

第二组形式处理的是“同一个系统在不同观测尺度上表现为固体还是流体”的问题。德博拉数来自流变学，由 Markus Reiner 在 1964 年提出，定义为材料的弛豫时间与观测时间尺度之比。\cite{reiner1964deborah}

\begin{equation}
\mathrm{De} = \frac{\tau_{\text{relax}}}{\tau_{\text{obs}}}
\end{equation}

$\tau_{\text{relax}}$ 是系统内部结构重新调整所需的时间，$\tau_{\text{obs}}$ 是观察者关心这件事的时间尺度，两者同量纲，因此 $\mathrm{De}$ 无量纲。$\mathrm{De} \ll 1$ 时，系统在观测尺度内完成调整，表现为可以流动的流体；$\mathrm{De} \gg 1$ 时，系统来不及调整，表现为固体的弹性响应。用一个具体例子说明它的用法：一项技术对就业的影响，如果以季度为观测尺度（$\tau_{\text{obs}}$ 小），则 $\mathrm{De}$ 很大，系统表现为刚性，短期内劳动者无法转岗；如果以二十年为观测尺度，则 $\mathrm{De}$ 变小，同一系统的流动性显现出来。两种判断都不错，它们描述的是同一个系统在不同尺度上的行为，而争论常常发生在双方使用了不同的 $\tau_{\text{obs}}$ 却没有说明。

若系统近似满足单弛豫时间的Maxwell响应，周期输入下的损耗模量在角频率与弛豫时间之积接近1时达到峰值。将这种判断用于组织，需要先识别相应的周期输入与恢复机制。季度和二十年的观察窗可以改变对调整快慢的描述，本身不会改变系统的实际吸收成本。治理若要降低失配，应改变响应能力、输入频率或缓冲结构；延长统计窗口只改变观察口径。

这个名字的来源值得一提，因为它点出了这个量的适用范围。Reiner 取自《士师记》中底波拉之歌的一句“山在耶和华面前流动”，用来说明同一物质在足够长的时间尺度上可以表现出流动，而在短尺度上表现出固体。这个出处说明德博拉数从一开始就是一个关于尺度对比的量，而不是关于材料类别的量；它不把系统分成刚性的和流动的两类，只回答“在当前的观测窗里，它来得及调整吗”。因此它在本书里的用法也应当照此理解：给某个社会技术系统贴上“刚性”或“有弹性”的标签是没有意义的，有意义的说法是指明观测窗。

这个量在本书中有一个直接的用处：它把“技术性失业是短期的还是长期的”这类争论还原为一个可计算的问题：$\tau_{\text{relax}}$（劳动力市场的重组时间）与 $\tau_{\text{obs}}$（政策关心的时间窗）的比值。争议因此不再关于立场，而关于两个可以被测量的时间常数。同样的还原可以用在其他争论上：关于某项技术“来不及治理”的判断，实际上是说 $\tau_{\text{relax}}$ 大于治理机制的响应时间；关于某项制度“很有韧性”的判断，实际上是说它的响应时间短于它需要吸收的扰动的周期。两种说法都可以被具体的测量替代，而一旦被替代，争论就从立场转到数据。

$\tau_{\text{relax}}$ 在社会系统里同样需要一个可测的代理量，而这个代理量往往就是一个物理的周期。职业结构的重组速度受制于培养周期，一个需要十年上下才能完成训练并独立执业的职业，其人员存量不可能在一两年内改变，因此 $\tau_{\text{relax}}$ 至少是培养周期的量级。设备和基础设施的更新速度受制于折旧年限与建设周期，制度安排的调整速度受制于立法与复议的周期，这些都可以从公开文件里读出具体年数。把这几项里最长的一项当作 $\tau_{\text{relax}}$ 的下界估计，是一个粗糙但可用的起点，因为它只用到可以直接查证的年限，不需要对行为做假设。粗略估计已经足以支撑一类判断：当政策关心的时间窗短于这个下界时，任何关于结构已经调整完毕的结论都不成立，无论数据看上去多么支持它。

观测窗的选择本身带有立场成分，这一点需要承认而不是回避。同一项变化，以季度为窗口会得到需要立即干预的结论，以代际为窗口会得到属于正常调整的结论，而两种窗口都可以被论证为合理。承认这一点不会使这个量失效，原因在于 $\tau_{\text{relax}}$ 与窗口无关：它由系统的结构决定，不由观察者决定，因此争论可以收窄为窗口选择的问题，而窗口选择是可以被公开讨论和检查的。可操作的规范是：给出结论时同时写明所用的窗口，并说明在这个窗口下被略过了什么。做到这一条，不同窗口下的判断可以并存而不互相取消，它们各自指明了在自己那段时间尺度上必须处理的问题。

用这个量的时候还有一点容易被忽略：弛豫时间不是一个系统只有一个，而是一个机制对应一个。劳动力市场的重组、供应链的重建、法规的修订、设备存量的更新，各有自己的时间常数，彼此可以相差很多，因此说某个系统的德博拉数是多少，必须连同机制一起说明。在实践里可以这样做：把关心的问题拆到具体机制上，为每一项给出自己的时间常数，再判断哪一项与观测窗最接近。与观测窗最接近的那一项，通常就是决定结论的那一项。同一件事在不同研究里会得到不同的判断，因为它们关心的机制不同。

\subsection{剪切增稠}

第三组形式处理的是“施加得越快，系统越硬”的现象。正文第四章讨论生物干预时使用了这一形式，其一般表达为幂律

\begin{equation}
\eta(\dot\gamma) = \eta_b + k\,\dot\gamma^{\,n},\qquad n > 0
\end{equation}

其中 $\dot\gamma$ 是剪切速率，量纲为时间的倒数，在这里对应干预施加的速率；$\eta_b$ 是零剪切黏滞，对应系统在缓慢干预下的响应，记作 $\eta_b$ 是为了与本附录前一节广义黏滞式里的基准黏滞系数 $\eta_0$ 区分开；$k$ 是增稠系数，$n$ 是幂指数。$n > 0$ 表示施加越快、等效黏滞越高，即系统以更强的抵抗回应更快的干预，$n$ 越大上升越陡。这个形式对生物系统特别适用，因为活系统有主动的代偿机制：干预速率超过代偿响应速率时，系统不会被动接受，而会启动反向调节，表现为表观阻力上升。

与它相对的是剪切变稀，即幂指数取负值或使用饱和形式的情形：施加越快，系统越容易流动。两种行为在流变学里都被研究得很清楚，区别在于系统的微观结构在剪切下是被拆散还是被组织起来。把这两个方向都写出来，可以避免一种常见的误读，即把“阻力随速率上升”当成所有复杂系统的普遍规律。真实的情形是按速率分段的：在低速区系统可能有足够时间重组，表现为剪切变稀；超过某个速率之后代偿跟不上，转变为剪切增稠。第四章讨论的干预强度阈值，对应的正是这个转折点。

按速率分段这件事有一个可以指认的机制：代偿需要时间，代偿机制的响应速率决定了转折点的位置。干预速率低于代偿速率时，系统来得及把扰动吸收掉，表现为阻力不随速率上升甚至下降；干预速率超过代偿速率时，代偿本身开始被扰动打乱，系统只能用更强的抵抗回应，表观黏滞转为上升。所以转折点不是某个固定的阈值，而是 $1/\tau_{\text{compensation}}$ 量级的速率，其中 $\tau_{\text{compensation}}$ 是该系统代偿环节的响应时间。这个读法把剪切变稀与剪切增稠统一成同一个系统在不同速率下的两种表现，而不是两类不同的系统，它与德博拉数处理尺度问题的思路一致：一个量比值决定看到哪一种行为。它同时给出一个可检验的推论：如果人为延长某个系统的代偿时间，例如削弱其调节能力，那么转折点应当向更低的干预速率移动，曲线的低速率段也会相应缩短。这个推论在细胞、组织和社会组织三个层次上都可以设计实验来检验，检验的结果决定这一节的形式还能不能继续外推。

流变学里这类行为之所以可以被当作定律使用，是因为它的机制被查清过。以高浓度颗粒悬浮液为例，剪切速率提高到一定程度后，颗粒来不及沿流线重新排列，被迫形成瞬时的聚集结构，流动阻力随之跳升，玉米淀粉浆在快速搅动下变硬就是这一机制的直接表现。这里的关键不在于现象引人注意，而在于现象的来源可以被独立验证：通过流变测量与散射实验能够确认聚集结构的存在，因此幂律中的指数与起始速率都可以被赋予机制上的解释。本书把这个形式移到生物与制度系统时，缺少的正是这一步独立验证，现有的依据是形状相似，而形状相似本身不足以支撑借用。所以这一节的地位必须写明：它给出的是一个待检验的外推，检验的标准是能否在目标系统里指出与聚集结构相对应的机制。

如果要在生物与制度系统里找这个形状的候选，判断的线索是：系统是否有一个需要时间的调节回路，以及这个回路是否会在输入过快时被自身打乱。抗生素使用强度上升与耐药率上升之间的关系、耕作强度上升与地力下降之间的关系、考核频次上升与被考核行为失真之间的关系，都属于这一形状的候选，它们的共同点是调节机制在高强度输入下反而成为扰动的一部分\cite{laxminarayan2013antibiotic}。列出这些候选不等于确认它们，每一例都需要单独的实证工作：测量起始强度、测量曲线的斜率、确认调节回路的存在。把候选与结论分开写，是为了避免一种常见的错误，即用一个领域里已经成立的机制，去为另一个领域里仅仅形状相似的曲线提供担保。

速率继续上升以后，系统可能饱和、失稳、损伤或堵塞，转折由具体机制决定。幂律只描述测量覆盖的区间。报告需要列出速率跨度、重复测量与转折迹象，并比较幂律、阈值或饱和形式在更宽区间的表现。高剪切率行为不能由拟合区间内的指数单独决定。

辨识幂律需要足够宽的输入区间和重复测量。基准项$\eta_b$不可忽略时，$\ln\eta$对$\ln\dot\gamma$的斜率不恒等于$n$，应拟合完整表达式，或说明在哪个区间可以忽略基准。还应比较阈值与饱和形式，并报告区间外转折。

三组形式的地位并不相同。指数法则和德博拉数是描述性的工作假设，其形式由机制的类比推得，参数需要实证确定；剪切增稠幂律在流变学中有扎实的实验基础，本书把它移到生物和制度系统上是一个有待检验的外推。把三者的地位分清，是为了避免一种常见的错误：把一个在一个领域内被验证过的形式，当作在所有领域内都成立的定律来使用。
